\par\noindent\textbf{(1)} 取模平方可知 \(\partial_z|A|^2=0\)，故每个时间点的功率保持 \(P(t)=P_0e^{-t^2/T_0^2}\)。传播解 \(A(z,t)=A(0,t)e^{i\phi_{\rm NL}(t)}\)，其中 \(\phi_{\rm NL}(t)=\gamma zP_0e^{-t^2/T_0^2}\)。这是纯相位演化。
\par\noindent\textbf{(2)} 复场瞬时频率偏移 \(\delta\omega(t)=-\partial_t\phi_{\rm NL}=2\gamma zP_0(t/T_0^2)e^{-t^2/T_0^2}\)。其关于 \(t\) 为奇函数：脉冲前沿和后沿产生相反频移，虽然时间强度包络本身未变。
\par\noindent\textbf{(3)} 对 \(f(t)=t e^{-t^2/T_0^2}\) 求导 \(f'=e^{-t^2/T_0^2}(1-2t^2/T_0^2)\)，极值在 \(t=\pm T_0/\sqrt2\)。最大频移模为 \(|\delta\omega|_{\max}=(\sqrt2\gamma zP_0/T_0)e^{-1/2}\)。
\par\noindent\textbf{(4)} 按Parseval定理 \(\int|A(t)|^2dt\) 不变，频域谱密度却因为时间相关相位而向宽频带重新分布。若加入二阶色散 \(i\beta_2A_{tt}/2\)，振幅与相位会耦合并产生脉冲压缩或展宽，不能继续沿用纯SPM时域强度不变结论。
